% Jiří Lebl, Basic Analysis v6.3; ch-real-nums.tex, lines 684–776
\subsection{Archimedean property}

As we have seen, there are plenty of real numbers in any interval.  But
there are also infinitely many rational numbers in any interval.  The
following is one of the fundamental facts about the real numbers.
The two parts of the next theorem are actually equivalent, even though it may
not seem like that at first sight.


\begin{thm} \label{thm:arch}
%FIXMEevillayouthack
\pagebreak[3]
\leavevmode
\begin{enumerate}[(i)]
\item \label{thm:arch:i} \emph{(\myindex{Archimedean property})}%
\footnote{Named after the Ancient Greek mathematician
\href{https://en.wikipedia.org/wiki/Archimedes}{Archimedes of Syracuse}
(c.\ 287 BC -- c.\ 212 BC)\@.
This property is Axiom V from Archimedes' \myquote{On the Sphere and
Cylinder} (225 BC)\@.}
If $x, y \in \R$ and
$x > 0$, then there exists an $n \in \N$ such that
\begin{equation*}
nx > y .
\end{equation*}
\item \label{thm:arch:ii} \emph{($\Q$ is dense in $\R$\index{density of
rational numbers})} If $x, y \in \R$ and
$x < y$, then there exists an $r \in \Q$ such that
$x < r < y$.
\end{enumerate}
\end{thm}

\begin{proof}
Let us prove \ref{thm:arch:i}.  Divide through by $x$.
Then \ref{thm:arch:i} says that for every real number $t\coloneqq \nicefrac{y}{x}$,
we can find $n \in \N$ such that $n > t$.  In other words,
\ref{thm:arch:i} says that $\N \subset \R$ is not bounded above.
Suppose for contradiction that $\N$ is bounded above.  Let $b \coloneqq \sup \N$.
The number $b-1$ cannot possibly be an upper bound for $\N$ as it is strictly
less than $b$ (the least upper bound).  Thus there exists an $m \in \N$ such that $m > b-1$.
Add one to obtain $m+1 > b$, contradicting $b$ being an
upper bound.

\begin{myfigureht}
\myincludepdft{figdensofQ}{%
A diagram of the horizontal real line with several tall vertical
ticks marked that are equally spaced with the distance
between them being 1 over n.  The second tick
from the left is marked as m minus 1 all over n.
After this tick there is the number x marked with a short tick.
The next tall tick is marked with m over n.
Then comes the number y and then a tall tick marked
m plus 1 all over n.  Then two more tall ticks are shown.}
\caption{Idea of the proof of the density of $\Q$: Find $n$ such that $y-x >
\nicefrac{1}{n}$, then take the least $m$ such that $\nicefrac{m}{n} > x$.\label{figdensofQ}}
\end{myfigureht}
Let us tackle \ref{thm:arch:ii}.
See \figureref{figdensofQ}
for a picture of the idea behind the proof.
First assume $x \geq 0$.
Note that $y-x > 0$.
By \ref{thm:arch:i}, there exists an $n \in \N$ such that
\begin{equation*}
n(y-x) > 1
\qquad \text{or} \qquad
y-x > \nicefrac{1}{n}.
\end{equation*}
Again by \ref{thm:arch:i}, the set
$A \coloneqq \{ k \in \N : k > nx \}$ is nonempty.  By the
\hyperlink{wop:link}{well ordering property}
of $\N$, $A$ has a least element $m$.
As $m \in A$, we have $m > nx$.
Divide through by $n$ to get $x < \nicefrac{m}{n}$.
As $m$ is the least
element of $A$, $m-1 \notin A$.
If $m > 1$, then $m-1 \in \N$, but $m-1 \notin A$ and so $m-1 \leq nx$.
If $m=1$,
then $m-1 = 0$, and $m-1 \leq nx$ still holds as $x \geq 0$.
In other words,
\begin{equation*}
m-1 \leq nx \qquad \text{or} \qquad m \leq nx+1 .
\end{equation*}
On the other hand,
from $n(y-x) > 1$ we obtain $ny > 1+nx$.
Hence $ny > 1+nx \geq m$, and therefore $y > \nicefrac{m}{n}$.
Putting everything together, we obtain $x < \nicefrac{m}{n} < y$.
So take $r = \nicefrac{m}{n}$.

Now assume $x < 0$.  If $y > 0$, then just take $r=0$.  If
$y \leq 0$, then $0 \leq -y < -x$, and we
find a rational $q$ such that $-y < q < -x$.  Then take $r = -q$.
\end{proof}



% Jiří Lebl, Basic Analysis v6.3; ch-seq-ser.tex, lines 2171–2259
\subsection{Bolzano--Weierstrass theorem}

While it is not true that every bounded sequence is convergent, the
Bolzano--Weierstrass theorem tells us that we can at least find a convergent
subsequence.
The version of Bolzano--Weierstrass 
we present in this section is the Bolzano--Weierstrass for
sequences of real numbers.

\begin{thm}[Bolzano--Weierstrass]\index{Bolzano--Weierstrass theorem}\label{thm:bwseq}
Suppose a sequence $\{ x_n \}_{n=1}^\infty$ of real numbers is bounded.
Then there exists a convergent subsequence $\{ x_{n_i} \}_{i=1}^\infty$.
\end{thm}

\begin{proof}
\thmref{subseqlimsupinf:thm} says that there exists
a subsequence whose limit is $\limsup_{n\to\infty} x_n$.
\end{proof}

The reader might complain right now that 
\thmref{subseqlimsupinf:thm} is strictly stronger than the
Bolzano--Weierstrass theorem as presented above.  That is true.
However, 
\thmref{subseqlimsupinf:thm} only applies to the real line, but
Bolzano--Weierstrass applies in more general contexts (that is, in $\R^n$)
with pretty much the exact same statement.

As the theorem is so important to analysis, we present an explicit
proof.
The idea of the following proof also generalizes to different contexts.

\begin{proof}[Alternate proof of Bolzano--Weierstrass]
As the sequence is bounded, there exist two numbers $a_1 < b_1$
such that $a_1 \leq x_n \leq b_1$ for all $n \in \N$.
We will define a subsequence $\{ x_{n_i} \}_{i=1}^\infty$ and two
sequences $\{ a_i \}_{i=1}^\infty$ and $\{ b_i \}_{i=1}^\infty$, such that
$\{ a_i \}_{i=1}^\infty$ is monotone increasing, $\{ b_i \}_{i=1}^\infty$ is monotone decreasing,
$a_i \leq x_{n_i} \leq b_i$ and such that $\lim_{i\to\infty} a_i =
\lim_{i\to\infty} b_i$.   That
$\{ x_{n_i} \}_{i=1}^\infty$ converges then follows by the \hyperref[squeeze:lemma]{squeeze lemma}.

We define the sequences inductively.  We will define the sequences so that
for all $i$, we have $a_i < b_i$,
and that $x_n \in [a_i,b_i]$ for infinitely many $n \in \N$.
We have already defined $a_1$ and $b_1$.  We take $n_1 \coloneqq 1$, that is
$x_{n_1} = x_1$.
Suppose that up to some $k \in \N$,
we have defined the subsequence $x_{n_1}, x_{n_2}, \ldots,
x_{n_k}$, and the sequences $a_1,a_2,\ldots,a_k$
and $b_1,b_2,\ldots,b_k$.
Let $y \coloneqq \frac{a_k+b_k}{2}$.
Clearly
$a_k < y < b_k$.  If there exist infinitely many $j \in \N$
such that $x_j \in [a_k,y]$, then set $a_{k+1} \coloneqq a_k$, $b_{k+1}
\coloneqq y$,
and pick $n_{k+1} > n_{k}$
such that $x_{n_{k+1}} \in [a_k,y]$.  If there are not infinitely many 
$j$ such that 
$x_j \in [a_k,y]$, then it must be true that there are infinitely many $j \in
\N$ such that 
$x_j \in [y,b_k]$.  In this case pick $a_{k+1} \coloneqq y$, $b_{k+1}
\coloneqq b_k$,
and pick $n_{k+1} > n_{k}$
such that $x_{n_{k+1}} \in [y,b_k]$.

We now have the sequences defined.  What is left to prove is that
$\lim_{i\to\infty} a_i = \lim_{i\to\infty} b_i$.  The limits exist as the sequences
are monotone.  In the construction,
$b_i - a_i$ is cut in half in each step.  Therefore,
$b_{i+1} - a_{i+1} = \frac{b_i-a_i}{2}$.  By
\hyperref[induction:thm]{induction},
\begin{equation*}
b_i - a_i = \frac{b_1-a_1}{2^{i-1}} .
\end{equation*}

Let $x \coloneqq \lim_{i\to\infty} a_i$.  As $\{ a_i \}_{i=1}^\infty$ is monotone,
\begin{equation*}
x = \sup \{ a_i : i \in \N \} .
\end{equation*}
Let $y \coloneqq \lim_{i\to\infty} b_i = \inf \{ b_i : i \in \N \}$.  Since $a_i < b_i$ for
all $i$, we have $x \leq y$.
As the sequences are monotone, we have (why?) that for all $i$,
\begin{equation*}
y-x \leq b_i-a_i = \frac{b_1-a_1}{2^{i-1}} .
\end{equation*}
Because $\frac{b_1-a_1}{2^{i-1}}$ is arbitrarily small and $y-x \geq 0$,
we have $y-x = 0$.  Finish by the \hyperref[squeeze:lemma]{squeeze lemma}.
\end{proof}



% Jiří Lebl, Basic Analysis v6.3; ch-seq-ser.tex, lines 2681–2879
\section{Cauchy sequences}
\label{sec:cauchy}

\sectionnotes{less than a lecture}

Often we wish to describe a certain number 
by a sequence that converges to it.  In this case, it is impossible to use the
number itself in the proof that the sequence converges.
It would be nice if we could check for convergence without knowing
the limit.

\begin{defn}
A sequence $\{ x_n \}_{n=1}^\infty$ is a \emph{\myindex{Cauchy sequence}}%
\footnote{%
Named after the French mathematician
\href{https://en.wikipedia.org/wiki/Cauchy}{Augustin-Louis Cauchy} (1789--1857).} if
for every $\epsilon > 0$ there exists an $M \in \N$ such that
for all $n \geq M$ and all $k \geq M$, we have
\begin{equation*}
\sabs{x_n - x_k} < \epsilon .
\end{equation*}
\end{defn}

Informally, being Cauchy means that the terms of the sequence are eventually
all arbitrarily close to each other.  We might expect such a sequence to be
convergent, and we would be correct due to $\R$ having the
\hyperref[defn:lub]{least-upper-bound property}.  Before we prove this fact,
we look at some examples.

\begin{example}
The sequence $\{ \nicefrac{1}{n} \}_{n=1}^\infty$ is a Cauchy sequence.

Proof:  Given $\epsilon > 0$, find $M$ such that
$M > \nicefrac{2}{\epsilon}$.  Then for $n,k \geq M$,
we have $\nicefrac{1}{n} < \nicefrac{\epsilon}{2}$
and
$\nicefrac{1}{k} < \nicefrac{\epsilon}{2}$.  Therefore, for $n, k \geq M$,
we have
\begin{equation*}
\abs{\frac{1}{n} - \frac{1}{k}}
\leq
\abs{\frac{1}{n}} + \abs{\frac{1}{k}}
< \frac{\epsilon}{2} + \frac{\epsilon}{2} = \epsilon.
\end{equation*}
\end{example}

\begin{example}
The sequence $\bigl\{ {(-1)}^n \bigr\}_{n=1}^\infty$ is not a Cauchy sequence.

Proof:
Given any $M \in \N$, take $n \geq M$ to be any even number, and
let $k \coloneqq n+1$.  Then
\begin{equation*}
\begin{split}
\babs{{(-1)}^n - {(-1)}^k}
& =
\babs{{(-1)}^n - {(-1)}^{n+1}}
\\
& =
\babs{1 - (-1)} = 2 .
\end{split}
\end{equation*}
Therefore, for any $\epsilon \leq 2$ the definition cannot be satisfied, and
the sequence is not Cauchy.
\end{example}

%\begin{example}
%The sequence $\{ \frac{n+1}{n} \}_{n=1}^\infty$ is a Cauchy sequence.
%
%Proof:  Given $\epsilon > 0$, find $M$ such that
%$M > \nicefrac{2}{\epsilon}$.  Then for $n,k \geq M$,
%we have $\nicefrac{1}{n} < \nicefrac{\epsilon}{2}$
%and
%$\nicefrac{1}{k} < \nicefrac{\epsilon}{2}$.  Therefore, for $n, k \geq M$,
%we have
%\begin{equation*}
%\begin{split}
%\abs{\frac{n+1}{n} - \frac{k+1}{k}}
%& =
%\abs{\frac{k(n+1)-n(k+1)}{nk}}
%\\
%& =
%\abs{\frac{kn+k-nk-n}{nk}}
%\\
%& =
%\abs{\frac{k-n}{nk}}
%\\
%& \leq
%\abs{\frac{k}{nk}}
%+
%\abs{\frac{-n}{nk}}
%\\
%& = \frac{1}{n} + \frac{1}{k}
%< \frac{\epsilon}{2} + \frac{\epsilon}{2} = \epsilon .
%\end{split}
%\end{equation*}
%\end{example}

\begin{prop}
If a sequence is Cauchy, then it is bounded.
\end{prop}

\begin{proof}
Suppose $\{ x_n \}_{n=1}^\infty$ is Cauchy.  Pick an $M$ such that for all
$n,k \geq M$, we have $\sabs{x_n-x_k} < 1$.  In particular, 
for all $n \geq M$,
\begin{equation*}
\sabs{x_n - x_M} < 1 .
\end{equation*}
By the reverse triangle inequality,
$\sabs{x_n} - \sabs{x_M} \leq \sabs{x_n - x_M} < 1$.  Hence, for $n \geq M$,
\begin{equation*}
\sabs{x_n} < 1 + \sabs{x_M}.
\end{equation*}
Let
\begin{equation*}
B \coloneqq \max \bigl\{ \sabs{x_1}, \sabs{x_2}, \ldots,
\sabs{x_{M-1}}, 1+ \sabs{x_M} \bigr\} .
\end{equation*}
Then $\sabs{x_n} \leq B$ for all $n \in \N$.
\end{proof}

\begin{thm}
A sequence of real numbers is Cauchy if and only if it converges.
\end{thm}

\begin{proof}
Suppose $\{ x_n \}_{n=1}^\infty$ converges to $x$,
and
let $\epsilon > 0$ be given.
Then there 
exists an $M$ such that for $n \geq M$,
\begin{equation*}
\sabs{x_n - x} < \frac{\epsilon}{2} .
\end{equation*}
Hence, for $n \geq M$ and $k \geq M$,
\begin{equation*}
\sabs{x_n - x_k} = 
\sabs{x_n - x + x - x_k}
\leq \sabs{x_n-x} + \sabs{x-x_k} < \frac{\epsilon}{2} + \frac{\epsilon}{2} =
\epsilon .
\end{equation*}

Alright, that direction was easy.  Now suppose $\{ x_n \}_{n=1}^\infty$ is Cauchy.
We have shown that $\{ x_n \}_{n=1}^\infty$ is bounded.
For a bounded sequence, liminf and limsup exist, and this is
where we use the
\hyperref[defn:lub]{least-upper-bound property}.
If we show that
\begin{equation*}
\liminf_{n\to \infty} x_n = \limsup_{n\to\infty} x_n ,
\end{equation*}
then $\{ x_n \}_{n=1}^\infty$ must be convergent by \propref{liminfsupconv:prop}.


Define $a \coloneqq \limsup_{n\to\infty} x_n$ and
$b \coloneqq \liminf_{n\to\infty} x_n$.
By \thmref{subseqlimsupinf:thm}, there exist subsequences
$\{ x_{n_i} \}_{i=1}^\infty$ and
$\{ x_{m_i} \}_{i=1}^\infty$ such that
\begin{equation*}
\lim_{i\to\infty} x_{n_i} = a
\qquad \text{and} \qquad
\lim_{i\to\infty} x_{m_i} = b.
\end{equation*}
Given an $\epsilon > 0$,
there exists an $M_1$ such that
$\sabs{x_{n_i} - a} < \nicefrac{\epsilon}{3}$ for all $i \geq M_1$
and an $M_2$ such that
$\sabs{x_{m_i} - b} < \nicefrac{\epsilon}{3}$ for all $i \geq M_2$.
There also exists an $M_3$
such that
$\sabs{x_n-x_k} < \nicefrac{\epsilon}{3}$
for all $n,k \geq M_3$.
Let $M \coloneqq \max \{ M_1, M_2, M_3 \}$.
If $i \geq M$, then $n_i \geq M$ and $m_i \geq M$.  Hence,
\begin{equation*}
\begin{split}
\sabs{a-b} & =
\sabs{a-x_{n_i}+x_{n_i}
-x_{m_i}+x_{m_i}
-b} \\
& \leq
\sabs{a-x_{n_i}}
+ \sabs{x_{n_i} -x_{m_i}}
+ \sabs{x_{m_i} -b} \\
& <
\frac{\epsilon}{3}
+
\frac{\epsilon}{3}
+
\frac{\epsilon}{3}
= \epsilon .
\end{split}
\end{equation*}
As $\sabs{a-b} < \epsilon$ for all $\epsilon > 0$, we have $a=b$ and 
the sequence converges.
\end{proof}



% Jiří Lebl, Basic Analysis v6.3; ch-contfunc.tex, lines 1381–1689
\section{Extreme and intermediate value theorems}
\label{sec:minmaxint}

%mbxINTROSUBSECTION

\sectionnotes{1.5 lectures}

Continuous functions on closed and bounded intervals
are quite well behaved.

\subsection{Min-max or extreme value theorem}

Recall that $f \colon [a,b] \to \R$ is
\emph{bounded}\index{bounded function}\index{function!bounded}
if there exists a $B \in \R$ such that
$\babs{f(x)} \leq B$ for all $x \in [a,b]$.
For a continuous function on a closed and bounded interval, we have the following lemma.

\begin{lemma}
A continuous function $f \colon [a,b] \to \R$ is bounded.
\end{lemma}

\begin{proof}
We prove the claim by contrapositive.  Suppose $f$ is not bounded.
Then for each
$n \in \N$, there is an $x_n \in [a,b]$ such that
\begin{equation*}
\babs{f(x_n)} \geq n .
\end{equation*}
The sequence $\{ x_n \}_{n=1}^\infty$ is bounded as $a \leq x_n \leq b$.
By the \hyperref[thm:bwseq]{Bolzano--Weierstrass theorem},
there is a convergent subsequence $\{ x_{n_i} \}_{i=1}^\infty$.
Let $x \coloneqq \lim_{i\to\infty} x_{n_i}$.
Since $a \leq x_{n_i} \leq b$ for all $i$, we have $a \leq x \leq b$.
The sequence $\bigl\{ f(x_{n_i}) \bigr\}_{i=1}^\infty$ is not bounded 
as 
$\babs{f(x_{n_i})} \geq n_i \geq i$.
Thus $f$ is not continuous at $x$ as
\begin{equation*}
f(x)
=
f\Bigl( \lim_{i\to\infty} x_{n_i} \Bigr) ,
\qquad \text{but} \qquad
\lim_{i\to\infty} f(x_{n_i}) \enspace \text{does not exist.} \qedhere
\end{equation*}
\end{proof}

Notice a key point of the proof.
Boundedness of $[a,b]$ allows us to use Bolzano--Weierstrass,
while the fact that it is closed gives us that the limit is back in $[a,b]$.
The technique is a common one: Find a sequence with a certain property,
then use Bolzano--Weierstrass to make such a sequence that also converges.

Recall from calculus that $f \colon S \to \R$ achieves an
\emph{\myindex{absolute minimum}}\index{minimum!absolute}%
\index{achieves absolute minimum}
at $c \in S$ if
\begin{equation*}
f(x) \geq f(c) \qquad \text{for all } x \in S.
\end{equation*}
On the other hand, $f$ achieves an 
\emph{\myindex{absolute maximum}}\index{maximum!absolute}%
\index{achieves absolute maximum}
at $c \in S$ if
\begin{equation*}
f(x) \leq f(c) \qquad \text{for all } x \in S.
\end{equation*}
If such a $c \in S$ exists, then we say
$f$ \emph{achieves an absolute minimum (resp.\ absolute maximum) on
$S$}, and we call
$f(c)$ the \emph{absolute minimum (resp.\ absolute maximum)}.
%See \figureref{fig:minmax}.
\begin{myfigureht}
\myincludegraphics{minmax}{%
A graph  of a function on the interval from a to b.
The graph is between two horizontal lines,
one marked as the absolute maximum of f which equals
f of c, and
one marked as the absolute minimum of f which equals
f of d.}
\caption{$f \colon [a,b] \to \R$ achieves an absolute maximum $f(c)$ at
$c$, and an absolute minimum $f(d)$ at $d$.\label{fig:minmax}}
\end{myfigureht}

If $S$ is a closed
and bounded interval, then a continuous $f$
is not just bounded, it must achieve an absolute minimum and an absolute
maximum on $S$.

\begin{thm}[Minimum-maximum theorem / Extreme value theorem]
\index{minimum-maximum theorem}%
\index{maximum-minimum theorem}%
\index{extreme value theorem}%
A continuous function $f \colon [a,b] \to \R$
achieves both an absolute minimum and an absolute maximum on $[a,b]$.
\end{thm}

Again, we remark that it is important that the domain of $f$ is a closed and bounded interval $[a,b]$.

\begin{proof}
The lemma says that $f$ is bounded, so
the set $f\bigl([a,b]\bigr) = \bigl\{ f(x) : x \in [a,b] \bigr\}$ has a supremum and an infimum.
There exist sequences
in the set $f\bigl([a,b]\bigr)$ that approach its supremum and its infimum.
That is, there are sequences
$\bigl\{ f(x_n) \bigr\}_{n=1}^\infty$ and $\bigl\{ f(y_n)
\bigr\}_{n=1}^\infty$, where $x_n$ and $y_n$ are in $[a,b]$,
such that
\begin{equation*}
\lim_{n\to\infty} f(x_n) = \inf f\bigl([a,b]\bigr) \qquad \text{and} \qquad
\lim_{n\to\infty} f(y_n) = \sup f\bigl([a,b]\bigr).
\end{equation*}
We are not done yet; we need to find where the minima and the maxima are.
The problem is that the sequences $\{ x_n \}_{n=1}^\infty$ and
$\{ y_n \}_{n=1}^\infty$ need not converge.
We know $\{ x_n \}_{n=1}^\infty$ and $\{ y_n \}_{n=1}^\infty$ are bounded
(their elements belong to a bounded interval $[a,b]$).
Apply the 
\hyperref[thm:bwseq]{Bolzano--Weierstrass theorem}
to find
convergent subsequences
$\{ x_{n_i} \}_{i=1}^\infty$ and 
$\{ y_{m_i} \}_{i=1}^\infty$.  Let
\begin{equation*}
x \coloneqq \lim_{i\to\infty} x_{n_i}
\qquad \text{and} \qquad
y \coloneqq \lim_{i\to\infty} y_{m_i}.
\end{equation*}
As $a \leq x_{n_i} \leq b$ for all $i$, we have $a \leq x \leq b$.
Similarly, $a \leq y \leq b$.  So $x$ and $y$ are in $[a,b]$.
For convergence sequences,
a limit of a subsequence is the same as the limit of the
sequence.  Also, we can take a limit past the continuous function $f$.  So,
\begin{equation*}
\inf f\bigl([a,b]\bigr) = \lim_{n\to\infty} f(x_n)
= \lim_{i\to\infty} f(x_{n_i}) = 
f \Bigl( \lim_{i\to\infty} x_{n_i} \Bigr) = f(x) .
\end{equation*}
Similarly,
\begin{equation*}
\sup f\bigl([a,b]\bigr) = \lim_{n\to\infty} f(y_n)
= \lim_{i\to\infty} f(y_{m_i}) = 
f \Bigl( \lim_{i\to\infty} y_{m_i} \Bigr) = f(y) .
\end{equation*}
Hence, $f$ achieves an absolute minimum at $x$ and
an absolute maximum at $y$.
\end{proof}

\begin{example}
The function $f(x) \coloneqq x^2+1$ defined on the interval $[-1,2]$ achieves a minimum
at $x=0$ when $f(0) = 1$.  It achieves a maximum at $x=2$ where $f(2) = 5$.
Do note that the domain of definition matters.  If we instead took the domain
to be $[-10,10]$, then $f$ would no longer have a maximum at $x=2$.
Instead,
the maximum would be achieved at $x=10$ and also at $x=-10$.
\end{example}

We show by examples that the different hypotheses of the theorem are
truly needed.

\begin{example}
The function $f \colon \R \to \R$ defined by $f(x) \coloneqq x$
achieves neither a minimum nor a maximum.  So it is important that
we are looking at a bounded interval.
\end{example}

\begin{example}
The function $f \colon (0,1) \to \R$ defined by $f(x) \coloneqq \nicefrac{1}{x}$
achieves neither a minimum, nor a maximum.
It is continuous, but $(0,1)$ is not closed.
The values of the function are
unbounded as we approach $0$.  Also as we approach $x=1$, the values of the
function approach $1$, but $f(x) > 1$ for all $x \in (0,1)$.  There is
no $x \in (0,1)$ such that $f(x) = 1$.  So it is important that
we are looking at a closed interval.
\end{example}

\begin{example}
Continuity is important.
Define $f \colon [0,1] \to \R$ by 
$f(x) \coloneqq \nicefrac{1}{x}$ for $x > 0$ and let $f(0) \coloneqq 0$.
The function does not achieve a maximum.  The domain $[0,1]$ is closed and
bounded, but the problem is that
the function is not continuous at 0.
\end{example}

\subsection{Bolzano's intermediate value theorem}

Bolzano's intermediate value theorem is one of the cornerstones of analysis.
It is sometimes only called the intermediate value theorem or just
Bolzano's theorem.  To prove Bolzano's theorem we prove the
following simpler lemma.

\begin{lemma} \label{IVT:lemma}
Let $f \colon [a,b] \to \R$ be a continuous function.
Suppose $f(a) < 0$ and $f(b) > 0$. 
Then there exists a number $c \in (a,b)$
such that $f(c) = 0$.
\end{lemma}

\begin{proof}
We define two sequences $\{ a_n \}_{n=1}^\infty$
and $\{ b_n \}_{n=1}^\infty$ inductively:
\begin{enumerate}[(i)]
\item Let $a_1 \coloneqq a$ and $b_1 \coloneqq b$.
\item If $f\left(\frac{a_n+b_n}{2}\right) \geq 0$, let $a_{n+1} \coloneqq a_n$ and
$b_{n+1} \coloneqq \frac{a_n+b_n}{2}$.
\item If $f\left(\frac{a_n+b_n}{2}\right) < 0$, let $a_{n+1} \coloneqq \frac{a_n+b_n}{2}$ and
$b_{n+1} \coloneqq b_n$.
\end{enumerate}
\begin{myfigureht}
\myincludegraphics{bisect}{%
Graph of a function that crosses the
x-axis at c going upwards.  The interval
a sub 1 to b sub 1 is marked and f is negative
at a sub 1 and positive at b sub 1.
Next interval is a sub 2 which is equal to a sub 1
and b sub 2 is in the middle of the previous interval.
Again the function is negative at a sub 2 and positive
at b sub 2.  We continue with the intervals
with f being negative on the left and positive
on the right.
The next interval is from a sub 3 which equals
a sub 2 and b sub 3 which
is in the middle of the previous interval.  
For the next interval, a sub 4 is in the middle and
b sub 4 is equal to b sub 3.
The next a sub 5 is in the middle again and b sub 5 
is equal to b sub 4.  The number c is
inside all of these intervals.}
\caption{Finding roots (bisection method).\label{bisectfig}}
\end{myfigureht}
See \figureref{bisectfig} for an example of the first five steps.
If $a_n < b_n$, then $a_n < \frac{a_n+b_n}{2} < b_n$.  So
$a_{n+1} < b_{n+1}$.
As $a_1 = a < b = b_1$,
\hyperref[induction:thm]{induction} gives that
$a_n < b_n$ for all $n$.
Furthermore, $a_n \leq a_{n+1}$ and 
$b_n \geq b_{n+1}$ for all $n$, that is, the sequences are monotone.
As $a_n < b_n \leq b_1 = b$ and 
$b_n > a_n \geq a_1 = a$ for all $n$,
the sequences are also bounded.  Therefore, the
sequences converge.
Let $c \coloneqq \lim_{n\to\infty} a_n$ and $d \coloneqq \lim_{n\to\infty} b_n$,
where also $a \leq c \leq d \leq b$.  We need
to show that $c=d$.
Notice
\begin{equation*}
b_{n+1} - a_{n+1} = \frac{b_n-a_n}{2}.
\end{equation*}
By \hyperref[induction:thm]{induction},
\begin{equation*}
b_n - a_n = \frac{b_1-a_1}{2^{n-1}} = 2^{1-n} (b-a) .
\end{equation*}
As $2^{1-n}(b-a)$ converges to zero, we take the limit as $n$ goes to
infinity to get
\begin{equation*}
d-c = \lim_{n\to\infty} (b_n - a_n) =
\lim_{n\to\infty} 2^{1-n} (b-a) = 0.
\avoidbreak
\end{equation*}
In other words, $c=d$.

By construction, for all $n$,
\begin{equation*}
f(a_n) < 0
\qquad \text{and} \qquad
f(b_n) \geq 0 .
\end{equation*}
Since
$\lim_{n\to\infty} a_n = \lim_{n\to\infty} b_n = c$
and $f$ is continuous at $c$, we may take 
limits in those inequalities:
\begin{equation*}
f(c) = \lim_{n\to\infty} f(a_n) \leq 0
\qquad \text{and} \qquad
f(c) = \lim_{n\to\infty} f(b_n) \geq 0 .
\end{equation*}
As $f(c) \geq 0$ and 
$f(c) \leq 0$, we conclude $f(c) = 0$.
Thus also $c \neq a$ and $c \neq b$, so
$a < c < b$.
\end{proof}

\begin{thm}[Bolzano's intermediate value theorem] \label{IVT:thm}
\index{Bolzano's theorem}
\index{Bolzano's intermediate value theorem}
\index{intermediate value theorem}
Let $f \colon [a,b] \to \R$ be a continuous function.
Suppose $y \in \R$ is such that $f(a) < y < f(b)$
or $f(a) > y > f(b)$.  Then there exists a $c \in (a,b)$
such that $f(c) = y$.
\end{thm}

The theorem says that a continuous function on a closed interval
achieves all the values between the values at the endpoints.

\begin{proof}
If $f(a) < y < f(b)$, then define $g(x) \coloneqq f(x)-y$.  Then 
$g(a) < 0$ and $g(b) > 0$, and we apply \lemmaref{IVT:lemma}
to $g$ to find $c$.  If $g(c) = 0$, then $f(c) = y$.

Similarly, if $f(a) > y > f(b)$, then define $g(x) \coloneqq y-f(x)$.
Again, $g(a) < 0$ and $g(b) > 0$, and we apply \lemmaref{IVT:lemma} to
find $c$.
As before, if $g(c) = 0$, then $f(c) = y$.
\end{proof}



% Jiří Lebl, Basic Analysis v6.3; ch-contfunc.tex, lines 2004–2169
\subsection{Uniform continuity}

We made a fuss of saying that the $\delta$ in the definition of
continuity depended on the point $c$.  There are situations when it is
advantageous to be able to pick a $\delta$ independent of any point, and so we
give a name to this concept.

\begin{defn}
Let $S \subset \R$, and let $f \colon S \to \R$ be a function.
Suppose for every $\epsilon > 0$ there exists a $\delta > 0$
such that whenever $x, c \in S$ and
$\sabs{x-c} < \delta$, then $\babs{f(x)-f(c)} < \epsilon$.
Then we say $f$ is \emph{\myindex{uniformly continuous}}.
\end{defn}

A uniformly continuous function must be continuous.
The only difference in the definitions
is that in uniform continuity,
for a given $\epsilon > 0$ we pick a $\delta > 0$ that
works for all $c \in S$.  That is, $\delta$ can no longer depend on $c$;
it only depends on $\epsilon$.  The domain of definition
of the function makes a difference now.  A function that is not uniformly
continuous on a larger set may be uniformly continuous when restricted to a
smaller set.
Note that $x$ and $c$ are not treated any differently
in this definition.

\begin{example}
$f \colon [0,1] \to \R$ defined by $f(x) \coloneqq x^2$ is uniformly continuous.

Proof: Note that $0 \leq x,c \leq 1$.  Then
\begin{equation*}
\sabs{x^2-c^2} = \sabs{x+c}\sabs{x-c}
\leq \bigl(\sabs{x}+\sabs{c}\bigr) \sabs{x-c}
\leq (1+1)\sabs{x-c} .
\end{equation*}
Therefore, given $\epsilon > 0$, let $\delta \coloneqq \nicefrac{\epsilon}{2}$.
If $\sabs{x-c} < \delta$, then $\sabs{x^2-c^2} \leq 2 \sabs{x-c} < \epsilon$.

\medskip

On the other hand, $g \colon \R \to \R$ defined by $g(x) \coloneqq x^2$ is not uniformly
continuous.

Proof: Suppose it is uniformly continuous.
Then for every $\epsilon > 0$,
there would exist a $\delta > 0$ such that
if $\sabs{x-c} < \delta$, then $\sabs{x^2 -c^2} < \epsilon$.
Take $x > 0$ and let
$c \coloneqq x+\nicefrac{\delta}{2}$.  Write
\begin{equation*}
\epsilon >
\sabs{x^2-c^2} = \sabs{x+c}\sabs{x-c}
=
(2x+\nicefrac{\delta}{2})\nicefrac{\delta}{2} 
\geq 
\delta x .
\end{equation*}
Therefore, $x < \nicefrac{\epsilon}{\delta}$ for all $x > 0$, which is a
contradiction.
\end{example}


\begin{example}
The function $f \colon (0,1) \to \R$ defined by $f(x) \coloneqq \nicefrac{1}{x}$ is not
uniformly continuous.

Proof: Given $\epsilon > 0$, the inequality $\epsilon >
\abs{\nicefrac{1}{x}-\nicefrac{1}{y}}$ holds if and only if
\begin{equation*}
\epsilon >
\abs{\nicefrac{1}{x}-\nicefrac{1}{y}}
=
\frac{\sabs{y-x}}{\sabs{xy}} 
=
\frac{\sabs{y-x}}{xy} ,
\end{equation*}
or
\begin{equation*}
\sabs{x-y} < xy \epsilon .
\end{equation*}
Suppose $\epsilon < 1$.
We wish to see if a 
small $\delta > 0$ would work.
If $x \in (0,1)$ and 
$y = x+\nicefrac{\delta}{2} \in (0,1)$,
then $\sabs{x-y} = \nicefrac{\delta}{2} < \delta$.
We plug $y$ into the inequality above to get $\nicefrac{\delta}{2} <
x \bigl( x+\nicefrac{\delta}{2} \bigr) \epsilon < x$.
If the definition of uniform continuity is satisfied, then
the inequality $\nicefrac{\delta}{2} < x$ would hold for all small $x > 0$.
But that implies $\delta \leq 0$.
Therefore, no single $\delta > 0$ works for all points.
\end{example}

The examples show that if $f$ is defined on an interval that is either not closed
or not bounded, then $f$ can be continuous, but not uniformly continuous.
For a closed and bounded interval $[a,b]$, we can, however,
make the following statement.

\begin{thm} \label{unifcont:thm}
Let $f \colon [a,b] \to \R$ be a continuous function.  Then $f$
is uniformly continuous.
\end{thm}

\begin{proof}
We prove the statement by contrapositive.
Suppose $f$ is not uniformly continuous.  We will prove
that there is some
$c \in [a,b]$ where $f$ is not continuous.  Let us negate
the definition of uniform continuouity:
There exists an $\epsilon > 0$
such that for every $\delta > 0$, there exist points $x, y$ in $[a,b]$ with
$\sabs{x-y} < \delta$ and $\babs{f(x)-f(y)} \geq \epsilon$.

So for the $\epsilon > 0$ above,
we find sequences $\{ x_n \}_{n=1}^\infty$ and $\{ y_n \}_{n=1}^\infty$ such that
$\sabs{x_n-y_n} < \nicefrac{1}{n}$ and such that $\babs{f(x_n)-f(y_n)} \geq
\epsilon$.  By
\hyperref[thm:bwseq]{Bolzano--Weierstrass},
there exists a convergent subsequence
$\{ x_{n_k} \}_{k=1}^\infty$.  Let $c \coloneqq \lim_{k\to\infty} x_{n_k}$.
As $a \leq x_{n_k} \leq b$ for all $k$, we have $a \leq c \leq b$.  Estimate
\begin{equation*}
\sabs{y_{n_k} - c} =
\sabs{y_{n_k} - x_{n_k} + x_{n_k} - c} \leq
\sabs{y_{n_k} - x_{n_k}}
+
\sabs{x_{n_k}-c}
<
\nicefrac{1}{n_k} 
+
\sabs{x_{n_k}-c} .
\end{equation*}
As $\nicefrac{1}{n_k}$ and $\sabs{x_{n_k}-c}$ both go to zero when
$k$ goes to infinity, $\{ y_{n_k} \}_{k=1}^\infty$ converges and the limit
is $c$.  We now show that $f$ is not continuous at $c$.
Estimate
\begin{equation*}
\begin{split}
\babs{f(x_{n_k}) - f(c)} & =
\babs{f(x_{n_k}) - f(y_{n_k}) + f(y_{n_k}) - f(c)} \\
& \geq
\babs{f(x_{n_k}) - f(y_{n_k})} - \babs{f(y_{n_k}) - f(c)} \\
& \geq
\epsilon - \babs{f(y_{n_k})-f(c)} .
\end{split}
\end{equation*}
Or in other words,
\begin{equation*}
\babs{f(x_{n_k})-f(c)} 
+
\babs{f(y_{n_k})-f(c)}  \geq
\epsilon .
\end{equation*}
At least one of the sequences $\bigl\{ f(x_{n_k}) \bigr\}_{k=1}^\infty$  or
$\bigl\{ f(y_{n_k}) \bigr\}_{k=1}^\infty$ cannot converge to $f(c)$.
Otherwise the
left-hand side of the inequality would go to zero while the right-hand side is positive.
Thus $f$ cannot be continuous at $c$.
\end{proof}

As before, note what is key in the proof: We can apply Bolzano--Weierstrass
because the interval $[a,b]$ is bounded, and the limit of the subsequence is
back in $[a,b]$ because the interval is closed.



% Jiří Lebl, Basic Analysis v6.3; ch-der.tex, lines 552–791
\subsection{Relative minima and maxima}

We previously talked about absolute maxima and minima.  These are the tallest peaks and
the lowest valleys in the entire mountain range.  What
about peaks of individual mountains and bottoms of individual valleys?
The derivative, being a local concept, is like walking around in a fog; it
cannot tell you if you are on the highest peak, but it can tell you whether you are
at the top of some peak.

\begin{defn}
Let $S \subset \R$ be a set and
let $f \colon S \to \R$ be a function.  The function $f$ is said to have
a \emph{\myindex{relative maximum}}\index{maximum!relative}
at $c \in S$ if there exists a $\delta>0$
such that for all $x \in S$ where $\sabs{x-c} < \delta$,
we have $f(x) \leq f(c)$.
The definition of
\emph{\myindex{relative minimum}}\index{minimum!relative}
is analogous.
\end{defn}

\begin{lemma}\label{relminmax:lemma}
Suppose $f \colon (a,b) \to \R$ is differentiable at $c \in (a,b)$,
and $f$ has
a relative minimum or a relative maximum at $c$.  Then
$f'(c) = 0$.
\end{lemma}

\begin{proof}
Suppose $c$ is a 
relative maximum of $f$.  That is, there is a $\delta > 0$ such
that for every $x \in (a,b)$ where
$\sabs{x-c} < \delta$, we have $f(x)-f(c) \leq 0$.
Consider the difference
quotient.  If $c < x < c+\delta$, then
\begin{equation*}
\frac{f(x)-f(c)}{x-c} \leq 0 ,
\end{equation*}
and if $c-\delta < y < c$, then
\begin{equation*}
\frac{f(y)-f(c)}{y-c} \geq 0 .
\end{equation*}
See \figureref{fig:critpt} for an illustration.
\begin{myfigureht}
\myincludegraphics{critpt}{%
A graph of a function f is shown.
Three points are marked on the horizontal axis, from left
to right, y, c, and x.  The function has a maximum
at the point c.  The points corresponding to x, c, and y are
marked on the graph.
Through the points corresponding
to y and c is a secant line that slopes upwards and is
marked with slope equals f of y minus f of c the whole thing
over quantity y-c is greater than or equal to 0.
Through the points corresponding
to c and x is a secant line that slopes downwards and is
marked with slope equals f of x minus f of c the whole thing
over quantity x-c is less than or equal to 0.}
\caption{Slopes of secants at a relative maximum.\label{fig:critpt}}
\end{myfigureht}

As $a < c < b$, there exist
sequences $\{ x_n \}_{n=1}^\infty$ and
$\{ y_n \}_{n=1}^\infty$ in $(a,b)$
such that $c < x_n < c+\delta$ and
$c-\delta < y_n < c$ for all $n \in \N$, and such that
$\lim_{n\to\infty} x_n = \lim_{n\to\infty} y_n = c$.
Since $f$
is differentiable at $c$,
\begin{equation*}
0 \geq \lim_{n\to\infty} \frac{f(x_n)-f(c)}{x_n-c} 
=
f'(c)
=
\lim_{n\to\infty} \frac{f(y_n)-f(c)}{y_n-c} \geq 0.
\end{equation*}
We are done with a maximum.
For a minimum, consider
the function $-f$.
\end{proof}

For a differentiable function, a point where 
$f'(c) = 0$ is called a \emph{\myindex{critical point}}.  When $f$ is not
differentiable at some points,
it is common to also say that $c$ is a critical point
if $f'(c)$ does not exist.
The theorem says that a relative minimum or maximum at an interior point
of an interval must be a critical point.
As you remember from calculus, one finds minima and maxima of a function
by finding all the critical points together with the endpoints of
the interval and simply checking at which of these points
the function is largest or smallest.

\subsection{Rolle's theorem}

Suppose a function has the same value at both endpoints of an interval.
Intuitively, it ought to attain a minimum or a maximum in the interior of the
interval.
Then at such a minimum or a maximum, the derivative should be zero.
See \figureref{rollefig} for the geometric idea.
This is the content of Rolle's theorem%
\footnote{Named after the French mathematician
\href{https://en.wikipedia.org/wiki/Michel_Rolle}{Michel Rolle}
(1652--1719).}.

\begin{myfigureht}
\myincludegraphics{rollefig}{%
A graph of a function is drawn on the interval from a to b.
The function is zero at both a and b.  The function goes up to
a maximum,
then down to a minimum, then back up.  The maximum of the function occurs
at a point marked c.  A horizontal tangent line is drawn on the
corresponding point on the graph.  Note that it is parallel with
the horizontal axis.}
\caption{Point where the tangent line is horizontal, that is $f'(c) =
0$.\label{rollefig}}
\end{myfigureht}

\begin{thm}[Rolle] \label{thm:rolle}
\index{Rolle's theorem}
Let $f \colon [a,b] \to \R$ be a continuous function
differentiable on $(a,b)$ such that $f(a) = f(b)$.
Then there exists a $c \in (a,b)$ such that $f'(c) = 0$.
\end{thm}

\begin{proof}
As $f$ is continuous on $[a,b]$, it attains an absolute minimum and an
absolute 
maximum in $[a,b]$.  We wish to apply \lemmaref{relminmax:lemma}, and
so we need to find some $c \in (a,b)$ where $f$ attains a minimum or a
maximum.
Write $K \coloneqq f(a) = f(b)$.
If there exists an $x$ such that $f(x) > K$, then the absolute
maximum is larger than $K$ and hence occurs at some $c \in (a,b)$, and
therefore $f'(c) = 0$.  On the other hand, if there exists an $x$
such that $f(x) < K$, then the absolute minimum occurs at some
$c \in (a,b)$, and so $f'(c) = 0$.  If there is no $x$ such that
$f(x) > K$ or
$f(x) < K$, then $f(x) = K$ for all $x$ and then
$f'(x) = 0$ for all $x \in [a,b]$, so any $c \in (a,b)$ works.
\end{proof}

It is absolutely necessary that the derivative exists for all $x
\in (a,b)$.  Consider the function $f(x) \coloneqq \sabs{x}$ on $[-1,1]$.
Clearly $f(-1) = f(1)$, but there is no point $c$ where $f'(c) = 0$.

\subsection{Mean value theorem}

We extend \hyperref[thm:rolle]{Rolle's theorem}
to functions that attain different
values at the endpoints.

\begin{thm}[Mean value theorem] \label{thm:mvt}
\index{mean value theorem}
Let $f \colon [a,b] \to \R$ be a continuous function
differentiable on $(a,b)$.  Then there exists a point $c \in (a,b)$
such that
\begin{equation*}
f(b)-f(a) = f'(c)(b-a) .
\end{equation*}
\end{thm}

For a geometric interpretation of the mean value theorem, see
\figureref{mvtfig}.  The idea is that the value $\frac{f(b)-f(a)}{b-a}$
is the slope of the line between the points $\bigl(a,f(a)\bigr)$
and $\bigl(b,f(b)\bigr)$.
Then $c$ is the point such that $f'(c) = \frac{f(b)-f(a)}{b-a}$, that 
is, the tangent line at the point $\bigl(c,f(c)\bigr)$ has the same slope as the
line between $\bigl(a,f(a)\bigr)$ and $\bigl(b,f(b)\bigr)$.
The name comes from the fact that the slope of the secant line
is the mean value of the derivative, so the average derivative is achieved
in the interior of the interval.

The theorem follows from \hyperref[thm:rolle]{Rolle's theorem}
by subtracting from $f$ the affine linear function
with the same values at $a$ and $b$ as $f$.
The graph of this affine linear function is the 
secant line---the straight line through
$\bigl(a,f(a)\bigr)$ and $\bigl(b,f(b)\bigr)$---and
its derivative is
$\frac{f(b)-f(a)}{b-a}$.
Then we are looking for a point where this
difference has zero derivative.

\begin{myfigureht}
\myincludegraphics{mvtfig}{%
A graph of a function is drawn on the interval from a to b.
The function looks similar as before but it is no longer
zero at the endpoints, the left endpoint is negative
and the right is positive.  A straight line through the endpoints
of the curve is drawn, that is, it is a line through the point
(a, f of a) and (b, f of b).  At the point corresponding to 
c on the horizontal axis the tangent line to the function is
parallel to the line through the endpoints.}
\caption{Graphical interpretation of the mean value theorem.\label{mvtfig}}
\end{myfigureht}


\begin{proof}
Define the
function $g \colon [a,b] \to \R$ by
\begin{equation*}
g(x) \coloneqq
f(x)-
\left( f(b)+\frac{f(b)-f(a)}{b-a}(x-b) \right)
=
f(x)-
f(b)-\frac{f(b)-f(a)}{b-a}(x-b) .
\end{equation*}
The function $g$ is differentiable on $(a,b)$,
continuous on $[a,b]$, such that $g(a) = 0$ and $g(b) = 0$.  Thus there exists
a
$c \in (a,b)$ such that $g'(c) = 0$, that is,
\begin{equation*}
0 = g'(c) = f'(c)-\frac{f(b)-f(a)}{b-a} .
\end{equation*}
In other words,
$f(b)-f(a) = f'(c)(b-a)$.
\end{proof}

The proof generalizes.  By considering
$g(x) \coloneqq
\bigl(f(x)-f(b)\bigr)
\bigl(\varphi(b)-\varphi(a)\bigr)
-
\bigl(f(b)-f(a)\bigr)
\bigl(\varphi(x)-\varphi(b)\bigr)$,
one can prove the following version.  We leave the proof as an exercise.

\begin{thm}[Cauchy's mean value theorem] \label{thm:cauchymvt}
\index{Cauchy's mean value theorem}
Let $f \colon [a,b] \to \R$ and $\varphi \colon [a,b] \to \R$ be continuous
functions
differentiable on $(a,b)$.  Then there exists a point $c \in (a,b)$
such that
\begin{equation*}
\bigl(f(b)-f(a)\bigr)\varphi'(c) = f'(c)\bigl(\varphi(b)-\varphi(a)\bigr) .
\end{equation*}
\end{thm}



% Jiří Lebl, Basic Analysis v6.3; ch-riemann.tex, lines 1–391
\chapter{The Riemann Integral} \label{int:chapter}

%%%%%%%%%%%%%%%%%%%%%%%%%%%%%%%%%%%%%%%%%%%%%%%%%%%%%%%%%%%%%%%%%%%%%%%%%%%%%%

\section{The Riemann integral}
\label{sec:rint}

%mbxINTROSUBSECTION

\sectionnotes{1.5 lectures}

An integral is a way to \myquote{sum} the values of a function.
There is sometimes confusion among students of
calculus between the \emph{integral} and the \emph{antiderivative}.
The integral is (informally) the area under the curve, nothing else.
That we can compute an antiderivative using the integral is a nontrivial
result we must prove.  
We will define the \emph{Riemann integral}%
\footnote{Named after the German mathematician
\href{https://en.wikipedia.org/wiki/Riemann}{Georg Friedrich Bernhard Riemann}
(1826--1866).}
using the Darboux integral%
\footnote{Named after the French mathematician
\href{https://en.wikipedia.org/wiki/Darboux}{Jean-Gaston Darboux} (1842--1917).},
an equivalent but technically simpler definition.

\subsection{Partitions and lower and upper integrals}

We want to integrate a bounded function defined on an interval $[a,b]$.
We first define two auxiliary integrals that are defined for all
bounded functions.  Only then can we talk about the Riemann integral and
the functions which it can integrate, the Riemann integrable functions.

\begin{defn}
A \emph{\myindex{partition}} $P$ of $[a,b]$ is
a finite set of numbers $\{ x_0,x_1,x_2,\ldots,x_n \}$ such that
\begin{equation*}
a = x_0 < x_1 < x_2 < \cdots < x_{n-1} < x_n = b .
\end{equation*}
We write
\begin{equation*}
\Delta x_i \coloneqq x_i - x_{i-1} .
\end{equation*}
Suppose $f \colon [a,b] \to \R$ is bounded and $P$ is a partition of
$[a,b]$.
Define
\begin{align*}
& m_i \coloneqq \inf \, \bigl\{ f(x) : x_{i-1} \leq x \leq x_i \bigr\} , &
& M_i \coloneqq \sup \, \bigl\{ f(x) : x_{i-1} \leq x \leq x_i \bigr\} , \\
& L(P,f) \coloneqq
\sum_{i=1}^n m_i \Delta x_i , &
& U(P,f) \coloneqq
\sum_{i=1}^n M_i \Delta x_i .
\avoidbreak
\end{align*}
We call $L(P,f)$ the \emph{\myindex{lower Darboux sum}} and
$U(P,f)$ the \emph{\myindex{upper Darboux sum}}\index{Darboux sum}.
\glsadd{not:lowerdarbouxsum}
\glsadd{not:upperdarbouxsum}
\end{defn}

The geometric idea of Darboux sums is indicated in
\figureref{darbouxfig}.  The lower sum is the area of the shaded
rectangles, and the upper sum is the area of the entire
rectangles, shaded plus unshaded parts.  The width of the $i$th rectangle is $\Delta x_i$,
the height of the shaded rectangle is $m_i$, and the height
of the entire rectangle is $M_i$.

\begin{myfigureht}
\myincludegraphics{darbouxfig}{%
A diagram of a graph of a function in a bold black line, all above the x
axis.
The domain is divided into 8 intervals and
the endpoints are marked, from left to right, x sub 0 through x sub 8.
The interval between x sub 4 and x sub 5 is marked as being of width
Delta x sub 5.
Above each interval are two rectangles.  For example,
there is a shaded rectangle of width Delta x sub 5
with the bottom side being exactly the interval from
x sub 4 to x sub 5 and of height m sub 5, which is the minimum value of the
function on the interval.  The rectangle is extended upwards unshaded until
M sub 5 which is the maximum value of the function on the interval.}
\caption{Sample Darboux sums.\label{darbouxfig}}
\end{myfigureht}

\begin{prop} \label{sumulbound:prop}
Let $f \colon [a,b] \to \R$ be a bounded function.  Let $m, M \in \R$ be 
such that for all $x \in [a,b]$, we have $m \leq f(x) \leq M$.  Then for every partition $P$
of $[a,b]$,
\begin{equation}
\label{sumulbound:eq}
m(b-a) \leq
L(P,f) \leq U(P,f)
\leq M(b-a) .
\end{equation}
\end{prop}

\begin{proof}
Let $P$ be a partition of $[a,b]$.  Note that for all $i$, we
have $m \leq m_i \leq M_i \leq M$.  We also have
$\sum_{i=1}^n \Delta x_i = (b-a)$.  Therefore,
\begin{multline*}
m(b-a) =
m \left( \sum_{i=1}^n \Delta x_i \right)
=
\sum_{i=1}^n m \Delta x_i
\leq
\sum_{i=1}^n m_i \Delta x_i 
\\
\leq
\sum_{i=1}^n M_i \Delta x_i
\leq
\sum_{i=1}^n M \Delta x_i 
=
M \left( \sum_{i=1}^n \Delta x_i \right)
=
M(b-a) .
\end{multline*}
Hence, we get \eqref{sumulbound:eq}.  In particular, the sets of lower and
upper sums are bounded sets.
\end{proof}

\begin{defn}
As the sets of lower and upper Darboux sums are bounded, we define
\begin{align*}
& \underline{\int_a^b} f(x)\,dx \coloneqq
\sup \, \bigl\{ L(P,f) : P \text{ a partition of } [a,b] \bigr\} , \\
& \overline{\int_a^b} f(x)\,dx \coloneqq
\inf \, \bigl\{ U(P,f) : P \text{ a partition of } [a,b] \bigr\} .
\end{align*}
We call $\underline{\int}$\glsadd{not:lowerdarboux}
the \emph{\myindex{lower Darboux integral}} and
$\overline{\int}$\glsadd{not:upperdarboux} the
\emph{\myindex{upper Darboux integral}}\index{Darboux integral}.
To avoid worrying about the variable of integration, 
we often simply write
\begin{equation*}
\underline{\int_a^b} f \coloneqq
\underline{\int_a^b} f(x)\,dx 
\qquad \text{and} \qquad
\overline{\int_a^b} f \coloneqq
\overline{\int_a^b} f(x)\,dx  .
\end{equation*}
\end{defn}

If integration is to make sense, then the lower and upper Darboux
integrals should be the same number, as we want a single number to call
\emph{the integral}.  However, these two integrals may differ for
some functions.

\begin{example} \label{example:dirichletfunc}
Take the \myindex{Dirichlet function}
$f \colon [0,1] \to \R$, where $f(x) \coloneqq 1$ if
$x \in \Q$ and $f(x) \coloneqq 0$ if $x \notin \Q$.  Then
\begin{equation*}
\underline{\int_0^1} f = 0 \qquad \text{and} \qquad
\overline{\int_0^1} f = 1 .
\end{equation*}
The reason is that for any partition $P$ and every $i$, we have 
$m_i = \inf \bigl\{ f(x) : x \in [x_{i-1},x_i] \bigr\} = 0$  and
$M_i = \sup \bigl\{ f(x) : x \in [x_{i-1},x_i] \bigr\} = 1$.  Thus
\begin{equation*}
L(P,f) = \sum_{i=1}^n 0 \cdot \Delta x_i = 0 , \quad \text{and} \quad
U(P,f) = \sum_{i=1}^n 1 \cdot \Delta x_i = \sum_{i=1}^n \Delta x_i = 1  .
\end{equation*}
\end{example}

\begin{remark}
The same definition of $\underline{\int_a^b} f$ and
$\overline{\int_a^b} f$
is used when $f$ is defined on a larger set $S$ such that
$[a,b] \subset S$.  In that case, we use the restriction of $f$ to $[a,b]$
and we must ensure that the restriction is bounded on $[a,b]$.
\end{remark}

To compute the integral, we often take a partition $P$ and make it finer.
That is, we cut intervals in the partition into yet smaller pieces.

\begin{defn}
Let $P = \{ x_0, x_1, \ldots, x_n \}$ and
$\widetilde{P} = \{ \widetilde{x}_0, \widetilde{x}_1, \ldots,
\widetilde{x}_{\ell} \}$ be
partitions of $[a,b]$.  We say $\widetilde{P}$ is a
\emph{refinement}\index{refinement of a partition} of $P$
if as sets $P \subset \widetilde{P}$.
\end{defn}

That is, $\widetilde{P}$ is a refinement of a partition if it contains all the
numbers in $P$ and perhaps some other numbers in between.  For example,
$\{ 0, 0.5, 1, 2 \}$ is a partition of $[0,2]$ and
$\{ 0, 0.2, 0.5, 1, 1.5, 1.75, 2 \}$ is a refinement.
The main reason for introducing refinements is the following proposition.

\begin{prop} \label{prop:refinement}
Let $f \colon [a,b] \to \R$ be a bounded function, and let $P$
be a partition of $[a,b]$.  Let $\widetilde{P}$ be a refinement of $P$.
Then
\begin{equation*}
L(P,f) \leq L(\widetilde{P},f) 
\qquad \text{and} \qquad
U(\widetilde{P},f) \leq U(P,f) .
\end{equation*}
\end{prop}

\begin{proof}
The tricky part of this proof is to get the notation correct.
Let $\widetilde{P} = \{ \widetilde{x}_0, \widetilde{x}_1, \ldots,
\widetilde{x}_{\ell} \}$ be
a refinement of 
$P = \{ x_0, x_1, \ldots, x_n \}$.  Then
$x_0 = \widetilde{x}_0$ and 
$x_n = \widetilde{x}_{\ell}$.  In fact, there are integers
$k_0 < k_1 < \cdots < k_n$ such that $x_i = \widetilde{x}_{k_i}$ for
$i=0,1,2,\ldots,n$.

Let $\Delta \widetilde{x}_q \coloneqq \widetilde{x}_q - \widetilde{x}_{q-1}$ for
$q=0,1,2,\ldots,\ell$.
See \figureref{fig:refinement}.
We get 
\begin{equation*}
\Delta x_i
=
x_i - x_{i-1} =
\widetilde{x}_{k_i} - \widetilde{x}_{k_{i-1}} =
\sum_{q=k_{i-1}+1}^{k_i} 
\widetilde{x}_{q} - \widetilde{x}_{q-1}
=
\sum_{q=k_{i-1}+1}^{k_i} \Delta \widetilde{x}_q .
\end{equation*}
\begin{myfigureht}
\myincludepdft{figrefinement}{%
A diagram of an interval between x sub quantity i minus 1 and x sub i of
width Delta x sub i.
It is divided into three smaller intervals.
The left hand endpoint x sub quantity i minus 1
is also marked as tilde x sub quantity q minus 3
and tilde x sub quantity k sub quantity i minus 1.
Similarly the right hand endpoint x sub i
is also marked as tilde x sub q and tilde x sub k sub i.
The two inside points are marked as
tilde x sub quantity q minus 2 and
tilde x sub quantity q minus 1.
The three intervals are of length
Delta tilde x sub quantity q minus 2,
Delta tilde x sub quantity q minus 1,
and
Delta tilde x sub q.}
\caption{Refinement of a subinterval.  Notice $\Delta x_i =
\Delta \widetilde{x}_{q-2} +
\Delta \widetilde{x}_{q-1} +
\Delta \widetilde{x}_{q}$,
and also
$k_{i-1}+1 = q-2$ and
$k_{i} = q$.\label{fig:refinement}}
\end{myfigureht}

Let $m_i$ be as before and correspond to the partition $P$.
Let $\widetilde{m}_q \coloneqq \inf \bigl\{ f(x) : \widetilde{x}_{q-1} \leq x \leq
\widetilde{x}_q \bigr\}$.
Now, $m_i \leq \widetilde{m}_q$ for $k_{i-1} < q \leq k_i$.  Therefore,
\begin{equation*}
m_i \Delta x_i
=
m_i \sum_{q=k_{i-1}+1}^{k_i} \Delta \widetilde{x}_q
=
\sum_{q=k_{i-1}+1}^{k_i} m_i \Delta \widetilde{x}_q
\leq
\sum_{q=k_{i-1}+1}^{k_i} \widetilde{m}_q \Delta \widetilde{x}_q .
\end{equation*}
So
\begin{equation*}
L(P,f) =
\sum_{i=1}^n m_i \Delta x_i
\leq
\sum_{i=1}^n \,
\sum_{q=k_{i-1}+1}^{k_i} \widetilde{m}_q \Delta \widetilde{x}_q
=
\sum_{q=1}^{\ell}
\widetilde{m}_q \Delta \widetilde{x}_q = L(\widetilde{P},f).
\end{equation*}

The proof of $U(\widetilde{P},f) \leq U(P,f)$ is left as an exercise.
\end{proof}

Armed with refinements, we prove the following.
The key point of this next proposition is that
the lower Darboux integral is less than or equal to the upper Darboux
integral.

\begin{prop} \label{intulbound:prop}
Let $f \colon [a,b] \to \R$ be a bounded function.  Let $m, M \in \R$ be 
such that for all $x \in [a,b]$, we have $m \leq f(x) \leq M$.  Then
\begin{equation}
\label{intulbound:eq}
m(b-a) \leq
\underline{\int_a^b} f \leq \overline{\int_a^b} f
\leq M(b-a) .
\end{equation}
\end{prop}

\begin{proof}
By \propref{sumulbound:prop}, for every partition $P$,
\begin{equation*}
m(b-a) \leq L(P,f) \leq U(P,f) \leq M(b-a).
\end{equation*}
The inequality
$m(b-a) \leq L(P,f)$ implies $m(b-a) \leq \underline{\int_a^b} f$.
The inequality
$U(P,f) \leq M(b-a)$ implies $\overline{\int_a^b} f \leq M(b-a)$.

The middle inequality in
\eqref{intulbound:eq} is the main point of the proposition.
Let $P_1, P_2$ be partitions of $[a,b]$.  Define 
$\widetilde{P} \coloneqq P_1 \cup P_2$.
The set $\widetilde{P}$ is a partition of $[a,b]$, which
is a refinement of $P_1$ and a refinement of $P_2$.
By \propref{prop:refinement},
$L(P_1,f) \leq L(\widetilde{P},f)$ and
$U(\widetilde{P},f) \leq U(P_2,f)$.  So
\begin{equation*}
L(P_1,f) \leq L(\widetilde{P},f) \leq U(\widetilde{P},f) \leq U(P_2,f) .
\end{equation*}
In other words, for two arbitrary partitions $P_1$ and $P_2$, we have
$L(P_1,f) \leq U(P_2,f)$.  
Recall \propref{infsupineq:prop}, and take the supremum and
infimum over all partitions:
\begin{multline*}
\underline{\int_a^b} f = 
\sup \, \bigl\{ L(P,f) : P \text{ a partition of } [a,b] \bigr\}
\\
\leq
\inf \, \bigl\{ U(P,f) : P \text{ a partition of } [a,b] \bigr\}
=
\overline{\int_a^b} f . \qedhere
\end{multline*}
\end{proof}

\subsection{Riemann integral}

We can finally define the Riemann integral.  However, the Riemann
integral is only defined on a certain class of functions, called the
Riemann integrable functions.

\begin{defn}
Let $f \colon [a,b] \to \R$ be a bounded function such that
\begin{equation*}
\underline{\int_a^b} f(x)\,dx = \overline{\int_a^b} f(x)\,dx .
\end{equation*}
Then $f$ is said to be \emph{\myindex{Riemann integrable}}.
The set of Riemann integrable functions on $[a,b]$ is denoted
by $\sR\bigl([a,b]\bigr)$.\glsadd{not:integrablefunc}
When $f \in \sR\bigl([a,b]\bigr)$, we define\glsadd{not:riemannint}
\begin{equation*}
\int_a^b f(x)\,dx \coloneqq 
\underline{\int_a^b} f(x)\,dx = \overline{\int_a^b} f(x)\,dx .
\end{equation*}
As before, we often write
\begin{equation*}
\int_a^b f \coloneqq \int_a^b f(x)\,dx.
\end{equation*}
The number $\int_a^b f$ is called the \emph{\myindex{Riemann integral}}
of $f$, or sometimes simply the \emph{integral} of $f$.
\end{defn}

By definition, a Riemann integrable function is bounded.
Appealing to \propref{intulbound:prop}, we immediately obtain
the following proposition.  See also \figureref{fig:integralminmax}.

\begin{prop} \label{intbound:prop}
Let $f \colon [a,b] \to \R$ be a Riemann integrable function.
Let $m, M \in \R$ be 
such that $m \leq f(x) \leq M$ for all $x \in [a,b]$.  Then
\begin{equation*}
m(b-a) \leq
\int_a^b f
\leq M(b-a) .
\end{equation*}
\end{prop}
\begin{myfigureht}
\myincludegraphics{integralminmax}{%
A diagram of a rectangle in the plane
that is horizontally between a and b, and vertically between the x axis
and capital M.  The rectangle is divided into the upper part
that is white and contains the graph of some function, and the lower
part that is shaded.  The dividing horizontal line is labeled as
lowercase m.}
\caption{The area under the curve is bounded from above by
the area of the entire rectangle, $M(b-a)$, and from below by
the area of the shaded part, $m(b-a)$.\label{fig:integralminmax}}
\end{myfigureht}



% Jiří Lebl, Basic Analysis v6.3; ch-riemann.tex, lines 1118–1195
\subsection{Continuous functions}

Let us show that continuous functions are Riemann integrable.
We can even allow some discontinuities.  We start with a
function continuous on the whole closed interval $[a,b]$.

\begin{lemma} \label{lemma:contint}
If $f \colon [a,b] \to \R$ is a continuous function,
then $f \in \sR\bigl([a,b]\bigr)$.
\end{lemma}

\begin{proof}
As $f$ is continuous on a closed bounded interval, it is
bounded and
uniformly continuous.
Given $\epsilon > 0$, find a $\delta > 0$ such that
$\sabs{x-y} < \delta$ implies $\babs{f(x)-f(y)} < \frac{\epsilon}{b-a}$.

Let $P = \{ x_0, x_1, \ldots, x_n \}$
be a partition of $[a,b]$ such that $\Delta x_i < \delta$ for all $i = 1,2,
\ldots, n$.  For example,
take $n$ such that $\frac{b-a}{n} < \delta$, and
let $x_i \coloneqq \frac{i}{n}(b-a) + a$.
Then for all $x, y \in [x_{i-1},x_i]$, we have 
$\sabs{x-y} \leq \Delta x_i < \delta$, and so
\begin{equation*}
f(x)-f(y) \leq \babs{f(x)-f(y)} < \frac{\epsilon}{b-a} .
\end{equation*}
As $f$ is continuous on $[x_{i-1},x_i]$, it attains a maximum and a minimum
on this interval.
Let $x$ be a point where $f$ attains the maximum and $y$ be a point
where $f$ attains the minimum.  Then $f(x) = M_i$
and $f(y) = m_i$ in the notation from the definition of the integral.
Therefore,
\begin{equation*}
M_i-m_i = f(x)-f(y) < 
\frac{\epsilon}{b-a} .
\end{equation*}
And so
\begin{equation*}
\begin{split}
\overline{\int_a^b} f - 
\underline{\int_a^b} f 
& \leq
U(P,f) - L(P,f)
\\
& =
\left(
\sum_{i=1}^n
M_i \Delta x_i
\right)
-
\left(
\sum_{i=1}^n
m_i \Delta x_i
\right)
\\
& =
\sum_{i=1}^n
(M_i-m_i) \Delta x_i
\\
& <
\frac{\epsilon}{b-a}
\sum_{i=1}^n
\Delta x_i
\\
& =
\frac{\epsilon}{b-a} (b-a)
= \epsilon .
\end{split}
\end{equation*}
As $\epsilon > 0$ was arbitrary,
\begin{equation*}
\overline{\int_a^b} f = \underline{\int_a^b} f ,
\end{equation*}
and $f$ is Riemann integrable on $[a,b]$.
\end{proof}



% Jiří Lebl, Basic Analysis v6.3; ch-riemann.tex, lines 1576–1780
\section{Fundamental theorem of calculus}
\label{sec:ftc}

%mbxINTROSUBSECTION

\sectionnotes{1.5 lectures}

In this section we discuss and prove the
\emph{\myindex{fundamental theorem of calculus}}.
The entirety of integral calculus is built upon this theorem,
ergo the name.
The theorem relates the seemingly unrelated concepts of integral and
derivative.  It tells us how to compute the antiderivative of a function
using the integral and vice versa.

\subsection{First form of the theorem}

\begin{thm} \label{thm:FTCv1}
Let $F \colon [a,b] \to \R$ be a continuous function, differentiable
on $(a,b)$.  Let $f \in \sR\bigl([a,b]\bigr)$ be such that $f(x) = F'(x)$ for $x \in
(a,b)$.  Then
\begin{equation*}
\int_a^b f = F(b)-F(a) .
\end{equation*}
\end{thm}

It is not hard to generalize the theorem to allow a finite number of points
in $[a,b]$ where $F$ is not differentiable, as long as it is continuous.
This generalization is left as an exercise.

\begin{proof}
Let $P = \{ x_0, x_1, \ldots, x_n \}$ be a partition of $[a,b]$.
For each interval $[x_{i-1},x_i]$, use the
\hyperref[thm:mvt]{mean value theorem} to find a
$c_i \in (x_{i-1},x_i)$ such that
\begin{equation*}
f(c_i) \Delta x_i = F'(c_i) (x_i - x_{i-1}) = F(x_i) - F(x_{i-1}) .
\end{equation*}
See \figureref{fig:fundthmfig}, and
note that the area of the
$i$th rectangle is
$F(x_{i})-F(x_{i-1})$,
and the total area of all three rectangles pictured is
$F(x_{i+1})-F(x_{i-2})$.
The idea is that taking smaller and smaller subintervals,
the total area of all these rectangles converges to the integral of $f$.
\begin{myfigureht}
\myincludegraphics{fundthmfig}{%
A diagram of a graph of a function in dark bold
marked as y equals f of x equals F prime of x
and three subintervals of the x coordinates.
The middle subinterval is labeled as going from
x sub quantity i minus 1 to x sub i and is of length Delta x sub i.
A point c sub i is marked inside this subinterval and a
dashed line goes up vertically until f of c sub i where it hits
the graph of f.  A shaded rectangle of this height and the subinterval
as the base is drawn and labeled with 'area equals f of c sub i times Delta
x sub i equals F of x sub i minus F of x sub quantity i minus 1'.
The other two subintervals are similar except on the left
we replace i with i minus 1 and on the right we replace i with
i plus 1.}
\caption{Mean value theorem on subintervals of a partition
approximating the area under the curve.\label{fig:fundthmfig}}
\end{myfigureht}

Using the notation from the definition of the integral,
$m_i \leq f(c_i) \leq M_i$, and multiplying by $\Delta x_i$ gets
\begin{equation*}
m_i \Delta x_i \leq F(x_i) - F(x_{i-1}) \leq M_i \Delta x_i .
\end{equation*}
We sum over $i = 1,2, \ldots, n$ to get
\begin{equation*}
\sum_{i=1}^n m_i \Delta x_i
\leq \sum_{i=1}^n \bigl(F(x_i) - F(x_{i-1}) \bigr)
\leq \sum_{i=1}^n M_i \Delta x_i .
\end{equation*}
In the middle sum, all the terms except the first and last cancel 
and we end up with $F(x_n)-F(x_0) = F(b)-F(a)$.  The sums on the left
and on the right are the lower and the upper sums, respectively.  So
\begin{equation*}
L(P,f) \leq F(b)-F(a) \leq U(P,f) .
\end{equation*}
We take the supremum of $L(P,f)$ over all partitions $P$ and the left inequality
yields 
\begin{equation*}
\underline{\int_a^b} f \leq F(b)-F(a) .
\end{equation*}
Similarly, taking
the infimum of $U(P,f)$ over all partitions $P$ yields
\begin{equation*}
F(b)-F(a) \leq \overline{\int_a^b} f .
\end{equation*}
As $f$ is Riemann integrable, we have
\begin{equation*}
\int_a^b f =
\underline{\int_a^b} f \leq F(b)-F(a) \leq \overline{\int_a^b} f
= \int_a^b f .
\end{equation*}
The inequalities must be equalities and we are done.
\end{proof}

The theorem is used to compute integrals.  Suppose we know that
the function $f(x)$ is a derivative of some other function $F(x)$,
then we can find an explicit expression for $\int_a^b f$. 

\begin{example}
To compute
\begin{equation*}
\int_0^1 x^2 \,dx ,
\end{equation*}
we notice $x^2$ is the derivative of $\frac{x^3}{3}$.
The fundamental theorem says
\begin{equation*}
\int_0^1 x^2 \,dx =
\frac{1^3}{3}
-
\frac{0^3}{3}
= \frac{1}{3}.
\end{equation*}
\end{example}

\subsection{Second form of the theorem}

The second form of the fundamental theorem gives us a way to solve
the differential equation $F'(x) = f(x)$, where $f$ is a known
function and we are trying to find an $F$ that satisfies the equation.

\begin{thm} \label{thm:FTCv2}
Let $f \colon [a,b] \to \R$ be a Riemann integrable function.  Define
\begin{equation*}
F(x) \coloneqq \int_a^x f .
\end{equation*}
First, $F$ is continuous on $[a,b]$.  Second,
if $f$ is continuous at $c \in [a,b]$, then $F$ is differentiable at $c$
and $F'(c) = f(c)$.
\end{thm}

\begin{proof}
As $f$ is bounded, there is an $M > 0$
such that $\babs{f(x)} \leq M$ for all $x \in [a,b]$.  Suppose $x,y \in [a,b]$
with $x > y$.  Then
\begin{equation*}
\babs{F(x)-F(y)} =
\abs{\int_a^x f - \int_a^y f}
=
\abs{\int_y^x f}
\leq
M\sabs{x-y} .
\end{equation*}
By symmetry, the same also holds if $x < y$.
So $F$ is Lipschitz continuous and hence continuous.

Now suppose $f$ is continuous at $c$.
Let $\epsilon > 0$ be given.  Let $\delta > 0$ be such that
for $x \in [a,b]$,
$\sabs{x-c} < \delta$ implies $\babs{f(x)-f(c)} < \epsilon$.
In particular,
for such $x$, we have
\begin{equation*}
f(c)-\epsilon < f(x) < f(c) + \epsilon.
\end{equation*}
Thus if $x > c$, then
\begin{equation*}
\bigl(f(c)-\epsilon\bigr) (x-c) \leq \int_c^x f \leq
\bigl(f(c) + \epsilon\bigr)(x-c).
\end{equation*}
When $c > x$, then the inequalities are reversed.  Therefore,
assuming $x \neq c$, we get
\begin{equation*}
f(c)-\epsilon
\leq
\frac{\int_c^{x} f}{x-c}
\leq
f(c)+\epsilon .
\end{equation*}
As 
\begin{equation*}
\frac{F(x)-F(c)}{x-c}
=
\frac{\int_a^{x} f - \int_a^{c} f}{x-c}
=
\frac{\int_c^{x} f}{x-c} ,
\end{equation*}
we have 
\begin{equation*}
\abs{\frac{F(x)-F(c)}{x-c} - f(c)} \leq \epsilon .
\end{equation*}
The result follows.  It is left to the reader to see why is it OK that we
just have a non-strict inequality.
\end{proof}

Of course, if $f$ is continuous on $[a,b]$, then it is automatically Riemann
integrable, $F$ is differentiable on all of $[a,b]$ and $F'(x) = f(x)$ for
all $x \in [a,b]$.

\begin{remark} \label{remark:fundthmbase}
The second form of the fundamental theorem of calculus still holds if
we let $d \in [a,b]$ and define
\begin{equation*}
F(x) \coloneqq \int_d^x f .
\end{equation*}
That is, we can use any point of $[a,b]$ as our base point.  The proof is
left as an exercise.
\end{remark}

